% Generated draft: requires mathematical and source review.
\documentclass[11pt]{article}
\usepackage[margin=25mm]{geometry}
\usepackage{fontspec}
\setmainfont{Latin Modern Roman}
\usepackage{amsmath,amssymb,mathtools,unicode-math}
\setmathfont{Latin Modern Math}
\usepackage{enumitem,needspace,xurl,hyperref}
\setlength{\parindent}{0pt}
\setlength{\parskip}{5pt}
\setlength{\emergencystretch}{3em}
\newcommand{\N}{\mathbb{N}}
\newcommand{\Z}{\mathbb{Z}}
\newcommand{\Q}{\mathbb{Q}}
\newcommand{\R}{\mathbb{R}}
\newcommand{\C}{\mathbb{C}}
\newcommand{\F}{\mathbb{F}}
\newcommand{\A}{\mathbb{A}}
\newcommand{\PP}{\mathbb{P}}
\newcommand{\E}{\mathbb{E}}
\newcommand{\eps}{\varepsilon}
\newcommand{\dd}{\,\mathrm{d}}
\begin{document}
\section*{UnsolvedMath ID: 20001876}
\textbf{Source problem code:} catalogue entry
\subsection{Definitions and mathematical statement}
Paternain's question asks how the sums of Betti numbers of length-filtered free-loop and fixed-endpoint path spaces grow, and whether this growth depends on the metric. The supplied text has damaged filtration notation. Following the aid, let \(M\) be closed, fix a coefficient field \(\mathbb k\), and use either the Sobolev free-loop space \(X=\Lambda M\) or path space \(X=\Omega(p,q)\). For a smooth Riemannian metric \(g\), set
\[
 L_g(\gamma)=\int_0^1|\dot\gamma|_g\,du,\qquad
 X_g(T)=\{\gamma\in X:L_g(\gamma)\leq T\},\qquad
 B_g(T)=\sum_{i\geq0}\dim_{\mathbb k}H_i(X_g(T);\mathbb k).
\]
Sums may be infinite unless finiteness hypotheses are supplied. This is a length filtration, not a silently substituted energy filtration. The source asks about the raw \(B_g\), which need not be monotone because sublevel classes can disappear under inclusion.

The source also attributes to [Gromov2001] a lower bound, for simply connected \(M\), of the intended form
\[
 B_g(CN)\geq\sum_{i=0}^N b_i(\Lambda M),
\]
for a suitable constant \(C\). The original superscripts and quantifiers are damaged; a complete citation and the precise constant, coefficient and filtration conventions are not supplied.

The paired aid treats persistent ranks instead. For \(m\geq0\), define
\[
 R_g^{\leq m}(s,t)=\sum_{i=0}^m
 \operatorname{rank}\bigl(H_i(X_g(s))\to H_i(X_g(t))\bigr),
\]
\[
 P_g^{\leq m}(T)=\sum_{i=0}^m
 \operatorname{rank}\bigl(H_i(X_g(T))\to H_i(X)\bigr).
\]
All homology uses \(\mathbb k\). These image ranks differ from raw sublevel Betti sums.

If \(A\geq1\) satisfies \(A^{-2}g\leq h\leq A^2g\) as quadratic forms, then \(A^{-1}L_g\leq L_h\leq AL_g\). Thus the two filtrations are multiplicatively \(A\)-interleaved. The reported two-scale consequence, for \(t\geq A^2s\), is
\[
 R_g^{\leq m}(s,t)\leq R_h^{\leq m}(As,t/A),
\]
with the same statement after exchanging \(g,h\). The image-to-total ranks obey
\[
 P_g^{\leq m}(T/A)\leq P_h^{\leq m}(T)\leq P_g^{\leq m}(AT).
\]
The analogous all-degree statements require meaningful finite sums. For a fixed total homology class, its minimax length, defined as the infimum of levels realizing that class in total homology, is also comparable by factors \(A^{-1},A\).

The aid concludes that the polynomial growth exponent and qualitative exponential-growth type of these image-rank functions are independent of the smooth metric, whenever the chosen growth quantities are defined. Numerical exponential rates need only be comparable after rescaling length; equality is not asserted. These comparisons do not determine the growth or metric dependence of the nonmonotone raw Betti sum asked about in the original question.
\subsection{Short English statement}
Uniform metric comparison rescales lengths and interleaves the loop and path filtrations. The supplied aid deduces metric-independent growth type for persistent-to-total image ranks, plus a two-scale persistent-rank inequality. It explicitly leaves the original raw sublevel Betti-sum question unresolved.
\subsection{Sources}
\begin{itemize}
\item Local source record: catalogue entry, catalogue entry; canonical \texttt{aim-geometry-notes.json}, \texttt{notes[213]}, source item 3.1, extracted Question 3.1.1 attributed to Paternain. Title: ``Metric comparison for filtered loop and path homology.''
\item AIM workshop ``Geodesics,'' supplied reference \url{https://aimath.org/pastworkshops/geodesicsproblems.pdf}. Original remarks and literature fields are empty. This draft uses the supplied record without accessing a PDF.
\item Original background reference [Gromov2001], concerning simply connected manifolds and a loop-space Betti lower bound; no full bibliographic entry is supplied. The persistence comparison theorem is a candidate contribution of research attempt 1.
\item Registry dated 14 August 2026: ``partially\_solved,'' ``PARTIAL-PROGRESS,'' statement status ``exact''; attempt partial, research-run status open, rigor label ``proved,'' novelty confidence low. The machine-generated, non-peer-reviewed aid requires independent verification before citation. The verification field introduces the supplied problem text without resolving its damaged notation.
\end{itemize}
\end{document}
